\honest{Pretending to be Wise}{Eliezer Yudkowsky}{2009}

I open with ``The hottest place in Hell is reserved for those who in time of crisis remain neutral,'' credit it to Dante, strike Dante out, and call John F. Kennedy a ``misquoter.'' \nb{Wikiquote traces the line to Henry Powell Spring in 1944; Kennedy popularized it as Dante's. Dante put those who were neither for nor against God near the mouth of Hell, not in its hottest place.}

You can be underconfident as well as overconfident. I single out one case: showing neutrality or suspended judgment ``in order to signal maturity, impartiality, or a superior vantage point.''

My parents answered theological questions, such as why ancient Egypt, which kept good records, has none of Jews ever being there, by saying they used to ask that sort of question too, ``but now I've grown out of it.'' \nb{The day before, in ``Against Maturity,'' Yudkowsky had explained a reply like this as a story the parents told themselves after giving up on such questions, not as a display for others.} A principal tells two children caught fighting, ``It doesn't matter who started the fight, it only matters who ends it.'' Of course it matters. If the principal does not know who started it, the principal should say so; let a parent punch the principal and see how far ``it doesn't matter who started it'' gets with a judge. But to adults a fight is only inconvenient, and it only matters that it end quickly. \nb{The second half of the principal's sentence gives the children a reason not to hit back. The post discusses only the first half.} Great Powers do the same, ``I believe,'' when they tell smaller groups to stop fighting: ``Oh, can't Israel and Hamas just get along?'' I name no such demand. \nb{The post appeared a month after the Gaza War of December 2008 and January 2009.}

I call this ``pretending to be Wise.'' Freire said that washing your hands of a conflict between the powerful and the powerless sides with the powerful. A world where the Great Powers only demand truces is good for aggressors and convenient for Great Powers, so ``part of this behavior can be chalked up to sheer selfishness on the part of the Wise.'' The rest is status: a principal who took sides would become ``a mere participant in the fray.'' So too the revered elder, a chief executive, a famous academic or the founder of a mailing list, who is respected as a judge on condition of almost never judging. \nb{The post shows that neutrality suits these people. It does not show that they are neutral because it suits them.}

Suspending judgment is sometimes rational. A commenter, Michael Rooney, noted a related error: beginning philosophy students, given reasons for doubt, conclude not that they should suspend judgment but that every view is as plausible as any other. But claiming that the evidence is balanced when it is not, between natural selection and intelligent design, say, is a fault. ``Neutrality is a definite judgment.'' It says the evidence supports only a neutral summing-up; it must predict something, it can be wrong, and it is as open to attack as any side. A call for compromise on abortion is a position like any other.

You may want to practice rationality first on topics less charged than abortion or the Israeli--Palestinian conflict. \nb{In ``Politics is the Mind-Killer'' Yudkowsky had advised using Louis XVI, not contemporary politics, to make a political point about rationality. This post's playground example would have done without abortion or Gaza.} Robin Hanson says divisive conversations are a limited resource, so you may spend yours on less common questions where you can ``pull the rope sideways'' and a little discussion does more good. Just do not pretend that neutral judgment, or declining to take part, makes you wiser than the people who take sides.
